\newpage
\part{Standard Configurations}

\section{Overview}

When configuring your AI agents (such as Claude Desktop, Cursor, Windsurf, or Gemini CLI), using official distributions pulled from canonical registries ensures a robust and capable setup. Depending on your toolchain, you can run these servers using Node's \lstinline[style=inline]|npx| or Python's \lstinline[style=inline]|uvx|.

Below is a summary of the most useful official MCP servers. Note that some functionalities (like filesystem and fetch) have implementations in both ecosystems, while others are unique to one.

\begin{paratable}
\footnotesize
\begin{tabularx}{\linewidth}{|>{\raggedright\arraybackslash}p{0.16\textwidth}|>{\raggedright\arraybackslash}X|>{\raggedright\arraybackslash}X|}
\hline
\textbf{Capability} & \textbf{NPX (Node.js)} & \textbf{UVX (Python)} \\
\hline
\textbf{Repository Control} &
\lstinline[style=inline]|@modelcontextprotocol/server-github| &
\lstinline[style=inline]|mcp-server-git| (local Git history) \\
\hline
\textbf{File System} &
\textit{N/A} &
\lstinline[style=inline]|mcp-server-filesystem| \\
\hline
\textbf{Fetch (HTML to MD)} &
\lstinline[style=inline]|@modelcontextprotocol/server-fetch| &
\lstinline[style=inline]|mcp-server-fetch| \\
\hline
\textbf{Web Search} &
\lstinline[style=inline]|@modelcontextprotocol/server-brave-search|\newline
\lstinline[style=inline]|exa-mcp-server| &
\textit{N/A} \\
\hline
\textbf{Database / Auditing} &
\lstinline[style=inline]|@modelcontextprotocol/server-postgres| &
\lstinline[style=inline]|mcp-server-sqlite| \\
\hline
\textbf{Reasoning \& Memory} &
\lstinline[style=inline]|@modelcontextprotocol/server-sequential-thinking|\newline
\lstinline[style=inline]|@modelcontextprotocol/server-memory| &
\textit{N/A} \\
\hline
\end{tabularx}
\end{paratable}
